% Chapter XIV commentary. Author: GPT-6 Astra (OpenAI, ChatGPT Pro mode).
% Source authority: attached edition, 7 October 2010, all eleven PDF leaves.

\begin{dcomment}{exposition}{p1}{What a square must remember}
Grothendieck wants a usable calculus of homotopy pullbacks on the category of homotopy types. The attraction is clear: four objects and four arrows are much easier to manipulate than a model category. The difficulty is already identified in his next paragraph. The comparison
\[
 \operatorname{dia}_{\square}:
 \mathcal M^{\square}(W^{\square})^{-1}
 \longrightarrow (\mathcal MW^{-1})^{\square}
\]
forgets part of a diagram. Here $\square$ has arrows $3\to1\to0$ and $3\to2\to0$, with the two composites equal; $\Phi\subset\square$ omits $3$. We use covariant diagrams, following the arrows actually drawn. An object on the left is a \emph{coherent square}; on the right it is an ordinary commutative square in the localized category. Calling an object coherent does not mean that its chosen representative has stopped commuting strictly. It means that the whole diagram was localized together.

Write $\operatorname{Cart}_{\mathrm{im}}$ for the replete image, on the right, of the homotopy-cartesian objects on the left. Membership says that \emph{some} coherent realization is cartesian. Proposition~1 will assert, incorrectly, that \emph{every} strict realization is then cartesian. An isomorphism of underlying squares need not lift to an isomorphism of coherent squares. Likewise, four arrows forming a cube after localization need not form a morphism of the chosen coherent squares. These are the two losses that the postface makes explicit.

The opening fibration formula belongs to the assumed theory of homotopy fibres. For an arbitrary Quillen model category, a safe version is that $Y_1\times_{Y_0}Y_2$ computes the homotopy pullback when $Y_1\to Y_0$ is a fibration and $Y_0,Y_2$ are fibrant; right properness permits the familiar version without these fibrancy assumptions \citep[Theorem~5.2.5 and the mixed Reedy argument of Lemma~5.2.6, dualized]{Hovey1999ModelCategories}. Thus the opening paragraph should not be detached from its hypotheses or used to diagnose a false formula in the transcription. The editors state the precise Quillen version below.

The commentary on Chapter~I, manuscript p.~7 (``Coherent squares and what the axioms remember''), and on Chapter~IV, manuscript p.~75 (``Constructing the loop square without choosing homotopies''), supplies the general background. The new issue here is reflection of a universal property through $\operatorname{dia}_{\square}$, and the chapter itself supplies a counterexample.
\end{dcomment}

\begin{dcomment}{context}{p2}{The earlier reason for keeping diagram categories}
The retreat to diagram categories on manuscript page~2 is part of a long-standing programme. In \emph{Pursuing Stacks}, \S69, Grothendieck relates it to the difficulty of formulating trace additivity and determinant multiplicativity using only a derived category. His proposed object retains the derived categories of diagram categories, including filtered complexes of varying lengths \citep[pp.~1--2 of the edited excerpt]{Grothendieck1983}. The letter to Thomason of 2~April~1991 returns to the mapping-cone problem and the need for this whole family \citep[\S2, pp.~5--6]{Grothendieck1991Letter}. These passages explain why an apparently modest square calculus leads back to derivators.

Heller's \emph{Homotopy Theories} already organizes homotopy Kan extensions on diagram categories \citep[Introduction, pp.~1--2]{Heller1988}. This is a historical parallel, not a warrant for all the proposed square rules. The point of the present fragment is precisely to test how much of that structure can survive on the single category $\mathcal H$.

Two different judgments in the source must therefore be distinguished. The ordinary square is too impoverished to carry the desired coherent universal property; nevertheless, some of the axioms imposed on such squares are too strong for their intended examples. There is no contradiction between these diagnoses. Nor should the editors' suggestion that Grothendieck stopped when he noticed the failure be converted into a documented account of his decision: the editors' postface proves the mathematical obstruction, not that biographical explanation.
\end{dcomment}

\begin{dcomment}{exposition}{p3}{Pasting, cancellation, and the cube}
The first three axioms try to isolate the part of pullback calculus that can be used without constructing a pullback. Cart~(2) has exactly the usual orientation: if the left vertical arrow is invertible, a square is distinguished precisely when the right vertical arrow is invertible. Cart~(1) supplies the transposed assertion. Cart~(3) then includes both composition and cancellation, with the left square held cartesian. Its failure for $\operatorname{Cart}_{\mathrm{im}}$ will be a failure of cancellation; the postface should not be read as disproving each implication separately.

We can fill in the cube argument entirely formally. Regard the cube as two horizontal squares $S_X,S_Y$ and name the opposite lateral faces
\[
 B=\left(\begin{matrix}X_1&\leftarrow&X_3\\\downarrow&&\downarrow\\Y_1&\leftarrow&Y_3\end{matrix}\right),
 \qquad
 F=\left(\begin{matrix}X_0&\leftarrow&X_2\\\downarrow&&\downarrow\\Y_0&\leftarrow&Y_2\end{matrix}\right).
\]
If $B$ and $S_Y$ are distinguished, their vertical composite $R$ is distinguished. The same rectangle is obtained by composing $S_X$ with $F$. If $F$ is distinguished, vertical cancellation implies that $S_X$ is distinguished. Vertical pasting is horizontal pasting after Cart~(1). This proves exactly the stronger three-face version in the first footnote; the other two lateral faces are unnecessary. An isomorphism of squares makes $B,F$ distinguished by Cart~(2), so the second footnote follows as well.

There is also an economy in the later list: Cart~(5)(b) is already a consequence of Cart~(1)--(3). For a cube between distinguished $S_X,S_Y$, suppose $p_0,p_1,p_2$ are isomorphisms. Then $F$ is distinguished. Composing $S_X$ and $F$ makes $R$ distinguished; cancelling $S_Y$ makes $B$ distinguished. Applying Cart~(2) to $B$ and $p_1$ forces $p_3$ to be an isomorphism. This deduction concerns the abstract axioms. The editors' two counterexamples will show why the intended class cannot satisfy them.
\end{dcomment}

\begin{dcomment}{exposition}{p5:sans être l'identité.}{Existence of a comparison is not a chosen comparison}
Cart~(4) supplies an object over each cospan. Cart~(5)(a) asks that a map of cospans extend to some map of these objects; it does not specify a choice compatible with composition. Cart~(5)(b), with its universal quantifier on the already chosen $p_3$, says that \emph{every} such extension over an isomorphism of cospans is invertible. Replacing ``every'' by ``there exists an invertible extension'' would be a different axiom.

For a fixed cospan $A\to C\leftarrow B$, consider the category of cones whose objects are triples $(T,a:T\to A,b:T\to B)$ satisfying $fa=gb$, and whose arrows preserve $a,b$. The first corollary says that the full subcategory consisting of distinguished completions is nonempty, that any two of its objects admit a comparison, and that every comparison is invertible. This is a connected groupoid, not a proof that it is a terminal category. Existence follows by Cart~(5)(a) with three identities; invertibility follows by (b). The source's nonuniqueness remark records the distinction between these conclusions. It supplies no separate example satisfying all five axioms with a nontrivial automorphism group.

For the intended homotopy class, the distinction is quite visible: the square with vertices $\Omega X,0,0,X$ and all edges zero permits every automorphism of $\Omega X$ over the fixed cospan. It also permits every \emph{noninvertible} endomorphism, including zero. Thus this example illustrates nonuniqueness but simultaneously invalidates (b); it cannot validate the literal axiom system. In an additive category satisfying the whole system, the development below proves that no such nonuniqueness is possible.

In the coherent repair the relevant base is itself a category of coherent cospans. Restriction from cartesian coherent squares to that category is an equivalence. This gives an honest functorial construction there, while forgetting coherence can again identify distinct maps. Neither a noncanonical choice of apex nor an arbitrary commuting cube restores the information that was forgotten.
\end{dcomment}

\begin{dcomment}{exposition}{p6:le carré h-cartésien n'est en général pas cartésien.}{Weak universality and the information in a homotopy}
The second corollary identifies the positive content of the existence axiom. A distinguished square is a \emph{weak pullback}: every compatible pair of maps from $T$ factors through its apex, without uniqueness. The constant square on $T$ is distinguished by Cart~(2), and applying Cart~(5)(a) to it gives exactly the surjection displayed in the text. Conversely, if every distinguished square is a weak pullback, apply its weak property to $T=X_3$ and the two composites $X_3\to X_i\to Y_i$. This supplies Cart~(5)(a). Thus, given the constant squares, the Hom-set assertion and that part of Cart~(5) are equivalent.

For an actual homotopy pullback the surjection has a geometric explanation. In an enhanced homotopy theory, put
\[
 M_i=\operatorname{Map}(T,Y_i).
 \qquad
 \operatorname{Map}(T,Y_3)\simeq M_1\mathbin{\times^h_{M_0}}M_2.
\]
A point of the homotopy pullback consists of two maps and a path between their composites. Equality in $\pi_0M_0$ ensures that such a path exists, which proves surjectivity after taking components. Different choices of path can give different components. For the loop square, taking $T=S^0$ in pointed spaces and $X=S^1$ gives the map $\pi_0\Omega S^1\cong\mathbb Z\to\{*\}$, a fully explicit failure of injectivity. This statement is about the intended homotopy class, not an example of all five literal axioms. Mapping spaces supply the universal property before $\pi_0$ is taken; compare \citet[\S1.2.13]{Lurie2009HTT}.

There is a useful rigidity argument whenever an ordinary pullback $P=A\times_C B$ already exists in $\mathcal H$. Let $Q$ be a distinguished completion. Its canonical map $q:Q\to P$ has a section $s:P\to Q$ by weak universality. The endomorphism $sq$ of $Q$ is over the identity cospan, so Cart~(5)(b) makes it invertible. But $(sq)^2=sq$; an invertible idempotent is the identity. Hence $q$ is an isomorphism, and $Q$ is an ordinary pullback. Conversely, Cart~(4) and invariance under isomorphisms show that every existing ordinary pullback square is distinguished. Consequently, if $\mathcal H$ has all ordinary pullbacks, the proposed class is forced to be precisely that class. This is the familiar minimality argument for weak universal objects; its weak-kernel form is explicitly recorded by \citet[\S1, ``Weak kernels'', p.~6 of the 2014 manuscript]{Krause2016FlatCovers}.
\end{dcomment}

\begin{dcomment}{later}{p7:mais c'est à peine un lien entre les deux structures}{Additivity, stability, and the two variables of Hom}
Grothendieck is looking for a common structure behind fibres and cofibres. Additivity alone cannot supply it. Take $\mathcal H=\mathrm{Vect}_k$, with ordinary pullback squares on the left and ordinary pushout squares on the right. Each structure satisfies the corresponding five axioms. For $V\ne0$, the square
\[
 \begin{tikzcd}[column sep=large]
 0\ar[d]&0\ar[l]\ar[d]\\
 V&0\ar[l]
 \end{tikzcd}
\]
is a pullback, since $0\times_V0=0$, but is not a pushout, since $0\amalg_00=0$. This refutes the additive speculation within the literal framework, without using any defective homotopy example.

For coherent diagrams, the corrected hypothesis is \emph{stability}. A pointed derivator is stable when its cartesian and cocartesian coherent squares coincide. The suspension--loop adjunction exists already in the pointed case; the equivalence of that adjunction characterizes stability \citetext{\citealp[Proposition~3.17]{Groth2013}; \citealp[Theorem~7.1, arXiv v2]{GrothPontoShulman2014}}. The corresponding higher-categorical statement appears in \citet[Proposition~1.1.3.4]{Lurie2017HA}. These theorems concern coherent squares, and do not rescue Cart~(3) on the ordinary homotopy category.

The Hom remark also has a precise scope. For fixed $T$, a derived mapping-space functor sends homotopy limits in its second variable to limits of spaces. In the first variable, mapping out of a homotopy colimit produces a homotopy limit. In a pointed enhancement this gives
\[
 \operatorname{Map}(\Sigma X,Y)
 \simeq\Omega\operatorname{Map}(X,Y)
 \simeq\operatorname{Map}(X,\Omega Y),
\]
where loops in the mapping space are based at the zero map; taking components gives the adjunction. Cisinski constructs the external derived Hom for model-derived derivators and proves its continuity \citep[Proposition~6.13, pp.~239--240]{Cisinski2003}. Internal Hom requires a specified closed monoidal structure and the appropriate derived adjunctions. Nothing on this page supplies such structure to an arbitrary category $\mathcal H$.
\end{dcomment}

\begin{dcomment}{development}{p7:mais c'est à peine un lien entre les deux structures}{Binary products force the literal structure to be ordinary}
The proposed axioms are meant to retain a weak universal property while allowing nonunique comparisons. The preceding rigidity argument shows that this freedom disappears wherever an ordinary pullback already exists. Pasting makes the restriction stronger: binary products alone force all the missing pullbacks to exist. This applies in particular to additive categories, without any assumption of kernels or of an abelian structure.

We assume Cart~(1)--(5) literally throughout this block. The weak-kernel minimality argument has the antecedent cited above. We checked the editors' postface, the coherent-square discussion of \citet[\S3]{Groth2013}, and \citet[\S1]{Krause2016FlatCovers}, together with targeted literature searches. We did not find the following deductions from these particular axioms there. The marking identifies proofs supplied by us, not a claim of priority.

\textbf{Pullback along a monomorphism.}\ours\quad
A distinguished square with a monomorphism among its two base arrows is an ordinary pullback.

\emph{Proof.} Write the square as the left square below, with $f$ monic. To test uniqueness of its lifting property, apply Cart~(4) to $P\xrightarrow qB\xleftarrow qP$ and obtain the right square:
\[
\begin{tikzcd}[column sep=large]
A\ar[d,"f"'] & P\ar[l,"p"']\ar[d,"q"]
 & R\ar[l,"r"']\ar[d,"s"]\\
C & B\ar[l,"g"] & P\ar[l,"q"] .
\end{tikzcd}
\]
Cart~(3) makes the outer rectangle distinguished. Its base cospan is $A\xrightarrow f C\xleftarrow{fp}P$, because $gq=fp$. Since $f$ is monic, this cospan has ordinary pullback $P$ with projections $p$ and $1_P$. The rigidity argument proved in the preceding exposition makes the canonical comparison $s:R\to P$ an isomorphism.

Weak universality of the right square supplies a diagonal $d:P\to R$ with $rd=sd=1_P$. Thus $d=s^{-1}$ and $r=s$. If $u,v:T\to P$ satisfy $qu=qv$, the same weak property gives $w:T\to R$ with $rw=u$ and $sw=v$; hence $u=v$. Therefore $q$ is monic. The original square is a weak pullback, and its factorization maps are unique because their composites with $q$ agree. It is an ordinary pullback. Cart~(1) treats the other base arrow.\hfill$\square$

\textbf{Collapse under binary products.}\ours\quad
If $\mathcal H$ has binary products and satisfies Cart~(1)--(5), it has all pullbacks, and its distinguished squares are exactly its ordinary pullback squares.

\emph{Proof.} Given $A\xrightarrow fC\xleftarrow gB$, form the cospan
\[
C\xrightarrow{\Delta}C\times C
  \xleftarrow{\,f\times g\,}A\times B.
\]
The diagonal is monic. A distinguished completion exists by Cart~(4), and the preceding result makes it an ordinary pullback. Its apex is $A\times_CB$: a pair $a:T\to A$, $b:T\to B$ with $fa=gb$ gives the unique map $(a,b)$ to $A\times B$ and the compatible map $fa$ to $C$, so the displayed pullback property gives a unique map to the apex. All pullbacks therefore exist. Rigidity shows that each distinguished square is ordinary; Cart~(4) and isomorphism invariance show that every ordinary pullback square is distinguished.\hfill$\square$

No terminal object was used. Additivity supplies the required binary products as biproducts, so the theorem is stronger than an additive statement. It does not classify categories lacking binary products, and we draw no general nonexistence conclusion in that setting.

\textbf{Pointed coincidence obstruction.}\ours\quad
If a pointed category carries both literal structures and their distinguished classes coincide, every object is isomorphic to zero.

\emph{Proof.} For any $X$, the cospan $0\to X\leftarrow0$ has ordinary pullback $0$. Cart~(4), rigidity, and isomorphism invariance designate the square with base $X$ and its other three vertices zero. Coincidence makes it distinguished cocartesian. Its upper span $0\leftarrow0\to0$ has ordinary pushout $0$; dual rigidity therefore gives $X\cong0$.\hfill$\square$

Even without coincidence, the literal loop construction is always zero, by the same pullback argument, and the literal suspension is zero if the right structure is present. Their adjunction would be the vacuous zero-functor adjunction. The meaningful suspension and loop functors in homotopy theory belong to the coherent setting anticipated on manuscript page~2.
\end{dcomment}

\begin{dcomment}{exposition}{p8:voir postface des éditeurs}{The precise assertion that fails in Proposition 1}
Grothendieck hopes to verify the axioms by transporting a homotopy-cartesian criterion through localization. The sufficient direction of Proposition~1 is the definition of $\operatorname{Cart}_{\mathrm{im}}$: a homotopy-cartesian strict square has a distinguished underlying image. The necessary direction would make homotopy cartesianness a property of the ordinary underlying square alone. The editors disprove precisely this reflection statement.

It matters that the obstruction occurs in Quillen model categories with entirely familiar diagram model structures. It cannot be removed by the promised availability of models for finite ordered indexing categories. The walking square and the walking cube are already finite posets; the example below lives in complexes over a field, where all their diagram model structures are available. An additional hypothesis could restrict the homotopy theory drastically, but Quillen presentability of diagrams does not establish the proposition.

The preliminary example, ordinary pullbacks in an ordinary category, remains valid. Symmetry, the isomorphism test, and pasting follow directly from the pullback universal property; every cospan has a completion by hypothesis; a map of cospans induces a unique map of pullbacks; an invertible base map induces an invertible comparison. These statements also prove the stronger uniqueness omitted from the proposed homotopical axioms. The preceding development gives a precise circumstance in which the general axiom system really does reduce to this example: the existence of binary products. It does not vindicate the intended model-category example.
\end{dcomment}

\begin{dcomment}{exposition}{p8:Il est équivalent de demander}{Postface: why one strict comparison cube suffices}
\emph{Editors' postface, printed p.~7.}
The editors replace the cospan before taking its limit. This makes the missing comparison map canonical, and explains the difference between a weak equivalence of whole squares and a mere isomorphism after applying $\operatorname{dia}_{\square}$.

Here is a verification with exactly their fibrancy hypotheses. Give $\Phi=(1\to0\leftarrow2)$ the Reedy degrees $\deg(2)=0$, $\deg(0)=1$, $\deg(1)=2$, with $2\to0$ direct and $1\to0$ inverse. A diagram $Y$ is Reedy fibrant precisely when $Y_0,Y_2$ are fibrant and $Y_1\to Y_0$ is a fibration. The constant-diagram functor is left Quillen: for a (trivial) cofibration its relative latching maps at $1,2$ are that map, and at $0$ the identity. Hence its right adjoint, the pullback functor, is right Quillen. This is the mixed Reedy method of \citet[Theorem~5.2.5 and the proof of Lemma~5.2.6]{Hovey1999ModelCategories}, dualized.

For a cospan $K$, choose a Reedy trivial cofibration $K\to RK$ with fibrant target. Any objectwise weak equivalence $K\to Y$ to a cospan with the printed fibrancy conditions extends to a map $RK\to Y$ by lifting against $Y\to *$. It is a weak equivalence by two-out-of-three. Applying the right Quillen pullback functor preserves it, since both diagrams are fibrant. Consequently the two comparison maps from the original fourth vertex are weak equivalences simultaneously. This proves ``some'' if and only if ``every''. Taking $Y_3=\lim Y$ makes the fourth component of the strict cube the unique induced cone map, so it also proves the single-cube criterion.

For completeness, right properness removes both objectwise fibrancy conditions. If $K_1\to K_0$ is a fibration and $K\to R$ is a fibrant replacement of the cospan, put $K'_1=R_1\times_{R_0}K_0$. Right properness and two-out-of-three make $K_1\to K'_1$ a weak equivalence between fibrant objects of the slice $\mathcal M/K_0$. Pullback to $K_2$ is right Quillen, so
\[
K_1\times_{K_0}K_2\longrightarrow
R_1\times_{R_0}K_2\longrightarrow R_1\times_{R_0}R_2
\]
are weak equivalences: the first by that slice argument, the second by right properness. This verifies the parenthesis in the source, without adding a properness assumption to the fibrant version.
\end{dcomment}

\begin{dcomment}{examples}{p8:Cette condition est très restrictive.}{Postface: a two-term complex and the missing filling}
\emph{Editors' postface, printed p.~8.}
The counterexample is especially transparent when every map can be written as a scalar. Let $k$ be a nonzero field and use unbounded cochain complexes, with quasi-isomorphisms as weak equivalences and degreewise surjections as fibrations \citep[Theorem~2.3.11, with degrees reversed]{Hovey1999ModelCategories}. Our shift convention is $(V[r])^n=V^{n+r}$, with differential $(-1)^r d_V$.

Take $X=k[0]$ and the acyclic complex $C=(k\xrightarrow{1}k)$ in degrees $0,1$. Let $p:C\to X$ be the identity in degree $0$ and zero in degree $1$. This is a surjective chain map. Its kernel $P$ is $k[-1]$, concentrated in degree $1$. Thus $0\xrightarrow{\sim}C\xrightarrow pX$ is the editors' factorization and
\[
\begin{tikzcd}[column sep=large]
C\ar[d,"p"'] & P\ar[l,"i"']\ar[d]\\
X & 0\ar[l]
\end{tikzcd}
\]
is both an ordinary and a homotopy pullback. The contraction $s^1:C^1\to C^0$ is the identity, and $d_Cs+sd_C=1_C$. In $D(k)$, therefore, $C\cong0$, and the underlying square is isomorphic to the square with apex $P$, base $X$, and all edges zero.

That latter square, taken \emph{strictly} in complexes, has comparison $s_3=0_P$ to the displayed fibrational replacement: commutativity requires $is_3=0$, and $i$ is injective. The comparison induces zero on $H^1(P)=k$, so it is not a quasi-isomorphism. This computes the dotted arrow in the postface and disproves the necessary direction of Proposition~1.

To locate the information lost, write down a homotopy-pullback model for $f:A\to Z\leftarrow B:g$:
\[
\begin{split}
H(f,g)^n&=A^n\oplus B^n\oplus Z^{n-1},\\
d(a,b,z)&=(d_Aa,d_Bb,fa-gb-d_Zz).
\end{split}
\]
The chain-map identities imply $d^2=0$. This is a shifted mapping cone of $f-g$, with the last coordinate negated relative to one common convention. A cone with maps $u:T\to A$, $v:T\to B$ and a specified degree $-1$ homotopy $h$ satisfying
\[
 d_Zh+hd_T=fu-gv
\]
has comparison $c_h=(u,v,h):T\to H(f,g)$. The displayed identity is exactly its chain-map condition. The coherent square is cartesian precisely when this comparison is a quasi-isomorphism.

For $A=B=0$, $Z=X=k[0]$, and $T=P=k[-1]$, the four edges are zero and $H(0,0)=P$. Every scalar $\lambda\in k$ specifies a filling $h_\lambda^1=\lambda$, and $c_{h_\lambda}$ is multiplication by $\lambda$ on $P$. It is a quasi-isomorphism exactly when $\lambda\ne0$. Its mapping cone is the two-term complex
\[
 k\xrightarrow{\lambda}k
 \quad\hbox{in degrees }0,1,
\]
which is acyclic for $\lambda\ne0$ and has $k$ in both cohomology degrees for $\lambda=0$. Thus identical vertices and edges support both a cartesian and a noncartesian coherent filling.

We can check the transport sign as well. Put $t=ps$, so $d_Xt+td_C=p$. The map
\[
 H(p,0)\longrightarrow X[-1],\qquad (c,z)\longmapsto z-tc
\]
is a chain map with acyclic kernel, the diagonal copy of $C$. Composing with $P\to H(p,0)$, $a\mapsto(i(a),0)$, gives $-1_P$. The strict pullback transported through this contraction has filling $h_{-1}$, whereas the strict zero-edge square has filling $h_0$. The sign depends on the displayed cone convention; its being nonzero does not. With the boundary and its identifications fixed, the classes of fillings are indexed by $k$, because $\underline{\operatorname{Hom}}(P,X)$ is concentrated in degree $-1$ and has no degree $-2$ boundaries. Allowing scalar automorphisms of the apex identifies all nonzero fillings, while the zero filling remains separate.

The editors' other interpretations agree with this computation. In pointed spaces, $X=S^1$ gives $\pi_0\Omega X=\mathbb Z$, so the zero self-map cannot be a weak equivalence. For a CW complex in general, the homotopy groups of its loop space detect all positive homotopy groups of the base component; Whitehead's theorem then makes exactly that component contractible \citep[Theorem~4.5 and \S4.3]{Hatcher2002}. Other components are unrestricted. For unbounded complexes in a Grothendieck abelian category, an acyclic middle term in $0\to P\to C\to X\to0$ gives $H^n(P)\cong H^{n-1}(X)$, so $P$ is acyclic exactly when $X$ is. The standard injective model structure supplies this setting \citep[Theorem~2.2 and Proposition~2.12 of the preprint]{Hovey2001Sheaves}; no assertion that acyclic complexes are chain-contractible is needed.
\end{dcomment}

\begin{dcomment}{status}{p8:dont les faces horizontales sont}{Postface: the two failed quantifiers}
\emph{Editors' postface, printed pp.~8--9.}
The two final counterexamples use the same object $P=\Omega X$ and the same underlying zero-edge square $Q$. They apply whenever $P$ is not isomorphic to zero, a condition realized by both worked examples above. In the horizontal pair, the left square and outer rectangle are $Q$, hence belong to $\operatorname{Cart}_{\mathrm{im}}$. In the right square the bottom edge is $1_0$ and the top edge is $0_P$. A homotopy pullback of an equivalence has an equivalence as opposite edge, so this right square cannot belong to $\operatorname{Cart}_{\mathrm{im}}$. This is exactly the implication ``I and III cartesian imply II cartesian'' that fails.

For Cart~(5)(b), use the endomorphism of $Q$ that is the identity at the three cospan vertices and zero at $P$. It is a commuting cube in the ordinary homotopy category, but its fourth component is not invertible. The positive existence statement in (a), and the failure of uniqueness already noticed in Corollary~2, should not be conflated with this false assertion about \emph{every} fourth component. As the formal deduction above shows, (b) would have followed from the full pasting axiom anyway.

The implication $0_P\in W\Rightarrow P\to0\in W$ used in the postface needs no cofibrancy assumption. Writing $r:P\to0$ and $s:0\to P$, the arrow $r$ is a retract, in the arrow category, of $sr=0_P$, using components $(1_P,s)$ and $(1_P,r)$. Weak equivalences are closed under retracts. The converse follows from two-out-of-three. Equivalently, an invertible zero endomorphism makes the object zero in the homotopy category. These are settled counterexamples, not conjectures awaiting a stronger model-category theorem.
\end{dcomment}

\begin{dcomment}{exposition}{p8:dont les faces horizontales sont}{Why composition nevertheless survives}
The failure of cancellation leaves a useful positive statement: for every Quillen model category, $\operatorname{Cart}_{\mathrm{im}}$ is closed under composition. Two adjacent cartesian witnesses can be made compatible along their shared edge. A left square and an outer rectangle, as in the counterexample, are different data. We give the rectification argument to keep this distinction precise.

For any model category $\mathcal A$, the arrow comparison
\[
U_{\mathcal A}:\operatorname{Ho}(\mathcal A^{[1]})
 \longrightarrow\operatorname{Ho}(\mathcal A)^{[1]}
\]
is full, essentially surjective, and conservative. The first two properties are the arrow-shape lifting property used in derivator strength \citep[Lemma~1.9]{Groth2013}. To verify fullness directly, take a Reedy-cofibrant source arrow $f:X_0\to X_1$ and a Reedy-fibrant target arrow $g:Y_0\to Y_1$. Thus $X_0$ is cofibrant, $f$ is a cofibration, and $Y_0,Y_1$ are fibrant. Represent a square in the homotopy category by maps $a_0,a_1$ and a right homotopy between $a_1f$ and $ga_0$. The endpoint map $\mathrm{ev}_0:\operatorname{Path}(Y_1)\to Y_1$ is a trivial fibration. Lifting against $f$ extends that homotopy from $X_0$ to $X_1$, replacing $a_1$ by a homotopic $a'_1$ with $a'_1f=ga_0$ strictly. This gives the required coherent-arrow morphism. Any arrow in the homotopy category has an actual representative between cofibrant-fibrant objects, proving essential surjectivity. Finally, pointwise weak equivalences and saturation imply conservativity. In particular, an isomorphism of underlying arrows lifts to an isomorphism of coherent arrows.

Now choose cartesian strict lifts of the two adjacent squares. Their common columns may differ, but the given underlying identifications specify an isomorphism between these columns in $\operatorname{Ho}(\mathcal M)^{[1]}$. Apply the preceding result for $\mathcal A=\mathcal M$ to lift it. The two squares can then be regarded as composable morphisms in $\operatorname{Ho}(\mathcal M^{[1]})$. Choose bifibrant representatives for the three column objects in the arrow model category, with a single shared middle representative, and represent the two morphisms by actual maps. They form one strict rectangle in $\mathcal M$.

Apply the arrow comparison a second time, now to $\mathcal A=\mathcal M^{[1]}$. Each new small square has an underlying arrow in $\operatorname{Ho}(\mathcal A)$ isomorphic to that of its original witness. Fullness lifts this isomorphism and conservativity makes the lift invertible in $\operatorname{Ho}(\mathcal M^{\square})$. Thus the two new squares are cartesian as coherent squares. Coherent pasting makes their outer rectangle cartesian, and its underlying square is isomorphic to the required composite. This proves the assertion without a combinatorial or stability assumption, and without lifting an arbitrary incoherent cube.

In the stable setting, \citet[Proposition~6.11, pp.~17--18 of the 7~June~2021 author version]{ChristensenFrankland2022} prove the corresponding pasting statement even for homotopy-cartesian squares in an arbitrary triangulated category. The valid remnant of the proposed calculus is therefore substantial: symmetry, the isomorphism test, completion of cospans, weak universality, and positive composition survive for the homotopy-image class. Cancellation and Cart~(5)(b) do not.
\end{dcomment}

\begin{dcomment}{later}{p8:dont les faces horizontales sont}{The coherent replacement of the square axioms}
\emph{After the editors' postface.}
The construction Grothendieck names on manuscript page~2 gives the appropriate replacement. Let $\mathbb D$ be a derivator and $j:\Phi\hookrightarrow\square$ the full inclusion omitting the prospective pullback vertex. A coherent square $Q\in\mathbb D(\square)$ is cartesian when the adjunction unit
\[
 Q\longrightarrow j_*j^*Q
\]
is invertible. Right Kan extension along $j$ is fully faithful; consequently
\[
 j^*:\mathbb D(\square)_{\mathrm{cart}}
       \simeq\mathbb D(\Phi)
\]
is an equivalence, with inverse $j_*$ \citep[Proposition~1.20 and Definition~3.9]{Groth2013}. A coherent cospan therefore has a cartesian completion, and a morphism of coherent cospans has a unique extension between these completions once their identifications are fixed. If its three components are isomorphisms, pointwise conservativity makes it an isomorphism in $\mathbb D(\Phi)$; the equivalence makes the extended square morphism invertible. This proves the correct fourth-vertex assertion. Its input is a morphism in the coherent category, not an arbitrary cube in $\mathbb D(e)$.

The full pasting-and-cancellation law is a statement about one coherent rectangle. To see the mechanism in a model or higher-category enhancement, write it as
\[
\begin{tikzcd}
A\ar[d]&B\ar[l]\ar[d]&C\ar[l]\ar[d]\\
D&E\ar[l]&F\ar[l].
\end{tikzcd}
\]
If the left square is cartesian, $B\simeq A\times^h_D E$. Associativity of the homotopy limit gives
\[
 B\times^h_E F\simeq
 (A\times^h_D E)\times^h_E F\simeq A\times^h_D F.
\]
Coherence identifies the comparison from $C$ to the first expression with its comparison to the last. One is an equivalence exactly when the other is. For derivators this is the right-Kan-extension pasting statement of \citet[Proposition~3.13]{Groth2013}; the higher-categorical proof is given, dually, by \citet[Lemma~4.4.2.1]{Lurie2009HTT}. The editors' rectangle loses the compatibility needed to make those two comparison maps agree.

In a stable enhancement there is a useful translation into the underlying triangulated category. A coherent cartesian square with projections $u:P\to A$, $v:P\to B$ and base maps $f:A\to C$, $g:B\to C$ supplies a distinguished triangle
\[
 P\xrightarrow{(u,-v)}A\oplus B
 \xrightarrow{(f,g)}C\xrightarrow{\delta}\Sigma P
\]
\citep[Theorem~6.1, arXiv v2]{GrothPontoShulman2014}. The connecting morphism is part of this conclusion. In the zero-edge loop example it is an isomorphism $X\to\Sigma P$, with the sign dictated by conventions. The cube with fourth component $0_P$ and base component $1_X$ cannot preserve that morphism, since it would require $\delta=(\Sigma0_P)\delta=0$. The triangle therefore diagnoses exactly why that cube is incoherent, even though its four visible faces commute in the homotopy category.

The relation with higher categories is a comparison of enhanced diagram theories. For a suitable model presentation, localizing diagram categories gives the homotopy categories of coherent diagrams; precise comparisons appear in \citet[Propositions~1.3.4.24--1.3.4.25]{Lurie2017HA}. It is not necessary here to identify every abstract derivator with an infinity category. The correct conclusion is narrower and sufficient: the familiar homotopy theories retain the diagram information on which the failed inference depended.
\end{dcomment}
